\subsection{定义一致结构的伪度量族的构造}

用伪度量族来定义一致结构的意义在于：所有的一致结构都能这样得到。即：

定理 1. 给定一致结构 \(\mathcal{U}\) （集合 \(X\) 上的），存在 \(X\) 上的一族伪度量，使得由这族伪度量定义的一致结构与 \(\mathcal{U}\) 相同。

对一致结构 U 的每个围域 V，递归地定义一列对称围域 \((\mathbf{U}_{n})\) ，使得 \(U_{1} \subset V\) 且 \(\dot{U}_{n+1}^{2} \subset U_{n}\) 对一切 \(n \geqslant 1\) 成立。序列 \((\mathbf{U}_{n})\) 是粗于 U 的某个一致结构 \(U_{v}\) 的一个围域基本系；此外，显然当 V 取遍 U 的围域滤子时，U 是诸一致结构 \(U_{v}\) 的上确界。因此定理 1 是下述命题的推论：

命题 2. 若 X 上的一致结构 U 具有可数的围域基本系，则存在 X 上的一个伪度量 f，使得 U 与由 f 定义的一致结构相同。

设 \((\mathbf{V}_n)\) 是 \(\mathcal{U}\) 的一个可数围域基本系。递归地定义 \((\mathbf{U}_n)\) （\(\mathcal{U}\) 的一列对称围域），使得 \(\mathbf{U}_1 \subset \mathbf{V}_1\) 且

\[
\mathbf {\dot {U}} _ {n + 1} ^ {3} \subset \mathbf {U} _ {n} \cap \mathbf {V} _ {n} \quad \text {   for   } \quad n \geqslant 1.
\]

显然 \((\mathbf{U}_n)\) 也是 \(\mathcal{U}\) 的一个围域基本系，且特别地有 \(\dot{\mathbf{U}}_{n+1}^3 \subset \mathbf{U}_n\) （对 \(n \geqslant 1\) ）。我们如下定义一个实值函数 \(g\) （在 \(\mathbf{X} \times \mathbf{X}\) 上）：\(g(x, y) = 0\) 若 \((x, y) \in \mathbf{U}_n\) 对一切 \(n\) 成立；\(g(x, y) = 2^{-k}\) 若 \((x, y) \in \mathbf{U}_n\) 对 \(1 \leqslant n \leqslant k\) 成立，但 \((x, y) \notin \mathbf{U}_{k+1}\) ；\(g(x, y) = 1\) 若 \((x, y) \notin \mathbf{U}_1\) 。函数 \(g\) 是对称的且是正的，并且 \(g(x, x) = 0\) 对一切 \(x \in \mathbf{X}\) 成立。令

\[
f (x, y) = \inf \sum_ {i = 0} ^ {p - 1} g \left(z _ {i}, z _ {i + 1}\right),
\]

其中下确界取遍由所有有限序列 \((z_{i})_{0\leqslant i\leqslant p}\) （ \(p\) 任意，满足 \(z_0 = x\) 且 \(z_{p} = y\) ）所成的集。我们将证明 \(f\) 是一个伪度量，并且满足不等式

\[
\frac {1}{2} g (x, y) \leqslant f (x, y) \leqslant g (x, y).\tag{2}
\]

由定义立即可得 \(f\) 是对称的、正的且满足三角不等式。同样显然有 \(f(x, y) \leqslant g(x, y)\) ，因而 \(f(x, x) = 0\) 对一切 \(x \in \mathbf{X}\) 成立，故 \(f\) 是一个伪度量。为证明不等式 (2) 的左半部分，我们对 \(p\) 用归纳法证明：对每个有限序列 \((z_i)_{0 \leqslant i \leqslant p}\) （由 \(p + 1\) 个 \(\mathbf{X}\) 的点组成，满足 \(z_0 = x\) 与 \(z_p = y\) ），有

\[
\sum_ {i = 0} ^ {p - 1} g (z _ {i}, \quad z _ {i + 1}) \geqslant \frac {\mathrm{I}}{2} g (x, y).\tag{3}
\]

当 \(p = \mathbf{i}\) 时这是明显的。令 \(a = \sum_{i=0}^{p-1} g(z_i, z_{i+1})\) ；若 \(a \geqslant \mathbf{i}/2\) ，则不等式 (3) 成立，因为 \(g(x, y) \leqslant \mathbf{i}\) 。于是设 \(a < \mathbf{i}/2\) ，并设 \(h\) 是使得下式成立的指标 \(q\) 中最大的一个

\[
\sum_ {i <   q} g (z _ {i}, z _ {i + 1}) \leqslant \frac {a}{2};
\]

；这时有 \(\sum_{i < h} g(z_i, z_{i+1}) \leqslant a/2\) 且 \(\sum_{i < h+1} g(z_i, z_{i+1}) > a/2\) ，从而

\[
\sum_ {i > h} g (z _ {i}, z _ {i + 1}) \leqslant \frac {a}{2}.
\]

由归纳假设有 \(g(x, z_h) \leqslant a\) 且 \(g(z_{h+1}, y) \leqslant a\) ；另一方面，显然有 \(g(z_h, z_{h+1}) \leqslant a\) 。设 \(k\) 是最小的整数 \(> 0\) ，使得 \(2^{-k} \leqslant a\) ；则 \(k \geqslant 2\) ，且 \((x, z_h) \in \mathbf{U}_k\) ，\((z_h, z_{h+1}) \in \mathbf{U}_k\) ，\((z_{h+1}, y) \in \mathbf{U}_k\) （由 \(g\) 的定义）；因而 \((x, y) \in \mathbf{U}_k \subset \mathbf{U}_{k-1}\) ，由此推出 \(g(x, y) \leqslant 2^{1-k} \leqslant 2a\) 。

于是不等式 (2) 得证；它们表明：对每个 \(a > 0\) ，集 \(\overline{f}^{1}([o, a])\) 包含 \(\mathbf{U}_{k}\) ，只要指标 \(k\) 满足 \(2^{-k} < a\) ；反之，每个 \(\mathbf{U}_{k}\) 包含集 \(\overline{f}^{1}([o, 2^{-k-1}])\) ；因而诸集 \(\overline{f}^{1}([o, a])\) 构成结构 \(\mathfrak{U}\) 的一个围域基本系。

注。一致结构 \(\mathfrak{U}\) （在 \(\mathbf{X}\) 上）由族 \(\Phi\) —— 由 \(\mathbf{X}\) 上在 \(\mathbf{X} \times \mathbf{X}\) 上一致连续的所有伪度量组成的族 —— 定义。因为显然由族 \(\Phi\) 定义的一致结构粗于 \(\mathfrak{U}\) ；反之，定理 1 表明存在 \(\Phi\) 的一个子族定义一致结构 \(\mathfrak{U}\) ，因而由 \(\Phi\) 定义的一致结构细于 \(\mathfrak{U}\) 。

